\documentclass[12pt,letterpaper]{article}
\usepackage[utf8]{inputenc}
\usepackage[spanish]{babel}
\usepackage{amsmath}
\usepackage{amsfonts}
\usepackage{amssymb}
\usepackage[margin=2.5cm]{geometry}

\begin{document}

\begin{center}
    \Large \textbf{Evaluación de Examen 2: Unidad II} \\
    \large Física para Ciencias de la Salud (005-1713) \\
    \vspace{0.5cm}
    \textbf{Estudiante:} María Meléndez \quad \textbf{C.I.:} 33.685.873
    \vspace{0.3cm} \\
    \large \textbf{Calificación Final: 9,5/10 puntos}
\end{center}

\vspace{0.5cm}

\section*{Análisis de Resultados}

\subsection*{1. Cinemática (Aceleración media)}
\textbf{Procedimiento y resultado:} Extrae y plantea la ecuación de aceleración media restando las velocidades ($0,6 - 0$) y dividiendo entre el tiempo ($0,15$). El cálculo aritmético es exacto ($4$); sin embargo, falla al indicar el análisis dimensional final, ya que encuadra el resultado como $4 \text{ m/s}$ (velocidad) en lugar de $\text{m/s}^2$. \\
\textbf{Veredicto:} \textbf{Parcialmente correcto} (Penalización leve por error en la unidad fundamental final) (1,5/2 pts).

\subsection*{2. Dinámica (Leyes de Newton)}
\textbf{Procedimiento y resultado:} Plantea la división matemática correcta derivando de la Segunda Ley de Newton ($150 \text{ N} / 25 \text{ kg}$). El cálculo es exacto ($6$). Aunque en el número final vuelve a escribir accidentalmente $6 \text{ m/s}$, la estudiante redacta una conclusión analítica donde indica textualmente ``una aceleración de 6 metros por segundo al cuadrado'', salvando así el concepto dimensional físico. \\
\textbf{Veredicto:} \textbf{Correcto} (2/2 pts).

\subsection*{3. Trabajo Mecánico}
\textbf{Procedimiento y resultado:} Plantea la multiplicación de la fuerza por la distancia ($400 \cdot 0,8$), indicando también el valor $1$ para representar el coseno del ángulo. El cálculo aritmético de $320 \text{ J}$ es perfectamente exacto. \\
\textbf{Veredicto:} \textbf{Correcto} (2/2 pts).

\subsection*{4. Energía Potencial}
\textbf{Procedimiento y resultado:} Extrae y plantea la ecuación de energía potencial gravitatoria multiplicando la masa, la gravedad y la altura ($1,2 \text{ kg} \cdot 9,8 \text{ m/s}^2 \cdot 1,5 \text{ m}$). El producto de los factores da el resultado verificado y correcto de $17,64 \text{ J}$. \\
\textbf{Veredicto:} \textbf{Correcto} (2/2 pts).

\subsection*{5. Potencia Mecánica}
\textbf{Procedimiento y resultado:} Extrae los datos y relaciona correctamente la fórmula dividiendo el trabajo entre el tiempo transcurrido ($P = \frac{W}{t}$). El desarrollo de la división de $75 \text{ J} / 60 \text{ s}$ es impecable, logrando la cifra correcta de $1,25 \text{ Watts}$. \\
\textbf{Veredicto:} \textbf{Correcto} (2/2 pts).

\vspace{0.5cm}
\section*{Comentarios Generales y Sugerencias}
¡Excelente examen! La estudiante demuestra un dominio sobresaliente de los conceptos de la Unidad II y un estupendo manejo de las operaciones aritméticas para llegar a los resultados correctos.
\begin{itemize}
    \item Se valora de forma sumamente positiva el hábito de **redactar una conclusión final explicativa** tras resolver los cálculos (visto en las preguntas 1 y 2). Esta práctica es excelente para interpretar físicamente los números en un contexto clínico o de laboratorio.
    \item Se aconseja, sin embargo, tener mayor precaución al transcribir las **unidades del Sistema Internacional**. Recordar siempre que en simbología matemática la aceleración lleva un exponente cuadrado ($\text{m/s}^2$).
    \item ¡Sigue con este mismo entusiasmo y nivel analítico en la próxima unidad!
\end{itemize}

\end{document}